\label{chap:83-09}

% Unidad U012. Biografía estructural completa de pi.

% La biografía de pi que seguía en esta residencia es exactamente la ya
% publicada en el capítulo 8, salvo el namespace de sus etiquetas. Se conserva
% en la fuente por su procedencia, pero no se imprime dos veces.
\iffalse

\section{Construction of the closure character and structural genealogical trajectories of \(\pi\)}
\label{p12-83-c9-s1:chap:biografia-pi-autonoma}

\subsection{Causal separation between seed, character, and evaluation}

The construction is divided into three strata:
\[
 \text{finite selection}
 \longrightarrow
 \text{construction of the closure character}
 \longrightarrow
 \text{Archimedean evaluation}.
\]
The decimal expansion does not intervene in the first stratum. The type chain
is
\begin{equation}
 104\,976
 \longrightarrow468\,\mathcal U_6
 \xrightarrow{\Psi}243\,\mathcal W_6
 \lhook\joinrel\longrightarrow\mathbb F_3^6,
 \qquad|\mathbb F_3^6|=729.
 \label{p12-83-c9-s1:eq:cadena-tipica-pi-autonoma}
\end{equation}
The cardinalities count, respectively, enriched initial conditions,
distinct emissions, visible ternary words, and elements
of the complete ambient fiber.

For \(U_6=(u_1,\ldots,u_6)\), the projection is
\[
 \Psi(U_6)=((-u_1)\bmod3,\ldots,(-u_6)\bmod3).
\]
The census of the \(43\) orbits of \(D_3\) proves that the closure
orbit is the only one whose six orientations possess at once:
\[
 \text{five lifts},\qquad
 \text{mass }1008,\qquad
 1008=432+4\cdot144.
\]
The gauge
\[
 \operatorname{diag}_c=6,
 \qquad
 \operatorname{card}=ES
\]
select a single representative:
\begin{equation}
 \boxed{w_6^\pi=010211.}
 \label{p12-83-c9-s1:eq:semilla-pi-autonoma}
\end{equation}
The selection usis catalogue, multiplicitiis, orbital action, and orientation;
it dois not consult to prefix of \(\pi\).

\subsection{Liftings \(L_0,L_1\) and 1st coinductive boundary}

On vectors row of \(\mathbb F_3^6\), let
\[
L_0=
\begin{pmatrix}
2&2&2&1&2&1\\
2&1&2&2&1&1\\
1&0&0&1&0&2\\
0&0&1&1&1&0\\
2&2&2&0&2&1\\
2&1&2&1&0&0
\end{pmatrix},
\quad
L_1=
\begin{pmatrix}
2&1&2&0&2&2\\
1&2&1&1&2&0\\
1&2&1&1&1&2\\
0&1&0&0&2&1\\
2&0&2&1&0&1\\
0&2&2&2&1&1
\end{pmatrix}.
\]
If \(b_1=w_6^\pi\), the recurrence of blocks is
\[
 b_{k+1}=b_kL_{\varepsilon_k},
 \qquad
 (\varepsilon_1,\varepsilon_2,\varepsilon_3,\varepsilon_4)=(0,0,1,1),
\]
and the accumulated word is formed by concatenation:
\[
 w_{6(k+1)}=w_{6k}\mathbin{\|}b_{k+1}.
\]
A word of length \(6k\) is not multiplied by a matrix \(6\times6\).

For the closing channel:
\[
\begin{array}{c|c}
\text{depth}&\text{new block}\\ \hline
6&010211\\
12&012222\\
18&010211\\
24&002111\\
30&110221
\end{array}
\]
and
\begin{equation}
 \boxed{
 w_{30}^{\pi}
 =010211|012222|010211|002111|110221.}
 \label{p12-83-c9-s1:eq:w30-pi-autonomo}
\end{equation}
The following boundary is
\[
 u_5^\pi=222220.
\]
Along with the other channels,
\[
 R_{36}=
 \begin{pmatrix}
 222220\\021222\\102011
 \end{pmatrix}.
\]
The symbol \(R_{36}\) records the sixth block of three histories and opens a
boundary; it is not a terminal.

The certified branch of twenty blocks begins
\begin{align*}
\pi:\;&010211|012222|010211|002111|110221|222220|111201|\\
&212121|200121|100100|101222|022212|012012|111210|\\
&121011|200220|120210|000101|022010|020111.
\end{align*}
For this reason, ni \(w_{30}\) ni \(R_{36}\) terminate the genealogy.

\subsection{Oriented decomposition of the microfiber of \texorpdfstring{\(\pi\)}{pi}: central region, \(4\) lateral regions, and multiplicity \texorpdfstring{\(432+4\cdot144=1008\)}{432+4x144=1008}}

The microfiber
\[
 \mathscr R_\pi^{\rm micro}
 :=\{U\in\mathcal U_6:\Psi(U)=010211\}
\]
contains exactly five regions:
\[
\begin{array}{c|c|c|c}
U_6&E_{\rm sig}&\operatorname{mult}&\text{role}\\ \hline
501|614|498|169|272|272&555555&432&\perp\\
810|923|870|169|272|272&339555&144&++\\
810|983|810|169|272|272&393555&144&++\\
870|923|810|169|272|272&933555&144&++\\
870|923|810|418|674|521&933717&144&--
\end{array}
\]
Consequently,
\begin{equation}
 \boxed{
 5=1_\perp+3_{++}+1_{--},
 \qquad
 1008=432+4\cdot144.}
 \label{p12-83-c9-s1:eq:descomposiciones-cinco-regiones-pi}
\end{equation}
The first identity classifies orientation functions. The second counts
provenance multiplicities. They do not count five indistinguishable copies.

The transversal region is
\[
 501|614|498|169|272|272,
 \qquad E_{\rm sig}=555555,
 \qquad\operatorname{mult}=432.
\]
The mutated return is
\[
 870|923|810|418|674|521,
 \qquad E_{\rm sig}=933717,
 \qquad\operatorname{mult}=144.
\]
Interchanging them would preserve the census \(3/1/1\) but falsify the complete
regional structure.

\subsection{Descent of a complete nonadic block to the phase transfer multigraph}

For \(U_6=(U_1,\ldots,U_6)\), let
\[
 A(U_6)=(U_1,U_2,U_3),
 \qquad
 B(U_6)=(U_4,U_5,U_6).
\]
Each emission defines an edge between \(A(U_6)\) and \(B(U_6)\). The raw
multigraph has \(96\) vertices, distributed into components of sizes
\[
 28,28,28,4,4,4.
\]
The internal phase is quotiented by
\[
 (a,b,c)\sim(b,c,a)\sim(c,a,b).
\]
The representative is the lexicographically minimal rotation. The quotient
graph has \(32\) vertices and two components of sizes \(28\) and \(4\).

The anchors are
\[
\begin{aligned}
U_{\rm APP}&=903|850|850|252|109|252,\\
U_e&=850|963|890|149|252|212,\\
U_\varphi&=830|943|830|109|252|252.
\end{aligned}
\]
We number the regions.
\[
\begin{aligned}
U_\pi^{(1)}&=810|923|870|169|272|272,\\
U_\pi^{(2)}&=810|983|810|169|272|272,\\
U_\pi^{(3)}&=870|923|810|169|272|272,\\
U_\pi^{(4)}&=870|923|810|418|674|521,\\
U_\pi^{(5)}&=501|614|498|169|272|272.
\end{aligned}
\]

\subsection{The \(5\) oriented closure cycles}

\subsubsection{Closure of \texorpdfstring{\(U_\pi^{(1)}\)}{U-pi 1}}
\begin{align*}
&252|109|252|903|850|850\to
903|850|850|252|109|252\to\\
&252|109|252|923|870|810\to
810|923|870|169|272|272\to\\
&272|169|272|963|890|850\to
850|963|890|149|252|212\to\\
&212|149|252|943|830|830\to
830|943|830|109|252|252\to
252|109|252|903|850|850.
\end{align*}

\subsubsection{Closure of \texorpdfstring{\(U_\pi^{(2)}\)}{U-pi 2}}
\begin{align*}
&903|850|850|252|109|252\to
272|169|272|903|850|850\to\\
&272|169|272|983|810|810\to
810|983|810|169|272|272\to\\
&272|169|272|963|890|850\to
850|963|890|149|252|212\to\\
&212|149|252|943|830|830\to
830|943|830|109|252|252\to
903|850|850|252|109|252.
\end{align*}

\subsubsection{Closure of \texorpdfstring{\(U_\pi^{(3)}\)}{U-pi 3}}
\begin{align*}
&903|850|850|252|109|252\to
292|189|232|903|850|850\to\\
&292|189|232|923|810|870\to
870|923|810|169|272|272\to\\
&272|169|272|963|890|850\to
850|963|890|149|252|212\to\\
&212|149|252|943|830|830\to
830|943|830|109|252|252\to
903|850|850|252|109|252.
\end{align*}

\subsubsection{Closure of \texorpdfstring{\(U_\pi^{(4)}\)}{U-pi 4}}
\begin{align*}
&903|850|850|252|109|252\to
272|169|272|903|850|850\to\\
&272|169|272|963|890|850\to
850|963|890|149|252|212\to\\
&212|149|252|923|810|870\to
870|923|810|418|674|521\to\\
&943|830|830|521|418|674\to
830|943|830|109|252|252\to
903|850|850|252|109|252.
\end{align*}

\subsubsection{Closure of \texorpdfstring{\(U_\pi^{(5)}\)}{U-pi 5}}
\begin{align*}
&252|109|252|903|850|850\to
903|850|850|252|109|252\to\\
&614|498|501|252|109|252\to
501|614|498|169|272|272\to\\
&272|169|272|963|890|850\to
850|963|890|149|252|212\to\\
&212|149|252|943|830|830\to
830|943|830|109|252|252\to
252|109|252|903|850|850.
\end{align*}

\begin{theorem}[Minimality of the Five Closures]
\label{p12-83-c9-s1:thm:minimalidad-cierres-pi-autonoma}
For each \(i\in\{1,\ldots,5\}\), the minimum length of a closed walk
that simultaneously traverses
\[
 U_\pi^{(i)},\quad U_e,\quad U_\varphi,\quad U_{\rm APP}
\]
in the quotient multigraph is eight.
\end{theorem}

\begin{proof}
Given \(i\), let
\[
 E_*=\{U_\pi^{(i)},U_e,U_\varphi,U_{\rm APP}\}.
\]
The finite space is considered.
\[
 \mathscr S_i=\{(v,M):v\in V(G),\ M\subseteq E_*\}.
\]
The first component records the vertex and the second the target edges
already traversed. If an edge \(e\) takes \(v\) to \(v'\), the transition is
\[
 (v,M)\longmapsto
 \bigl(v',M\cup(E_*\cap\{e\})\bigr).
\]
Breadth-first search enumerates lengths in increasing order. The first
return to \((v,E_*)\) occurs at length eight for all five regions.
The printed walks attain the bound; the optimality of the search excludes
a shorter length.
\end{proof}

The closure function is encoded by
\[
 \mathfrak C_\pi=
 \bigl(\mathscr R_\pi^{\rm micro},
 w_\pi,\rho_\pi,m_\pi,\ell_{\rm cl}\bigr),
\]
\[
 \rho_\pi=(\perp,++ ,++ ,++ ,--),
 \qquad
 m_\pi=(432,144,144,144,144),
 \qquad
 \ell_{\rm cl}=8.
\]

\subsection{Incidence operator and oriented half-turn}

The incidence shadow is
\[
 A_W=
 \begin{pmatrix}
 0&1&1&1&1&1\\
 1&0&1&2&2&1\\
 1&1&0&1&2&2\\
 1&2&1&0&1&2\\
 1&2&2&1&0&1\\
 1&1&2&2&1&0
 \end{pmatrix}.
\]
Multiplication in \(\mathbb F_3\) gives
\begin{equation}
 \boxed{A_W^2=2I_6=-I_6.}
 \label{p12-83-c9-s1:eq:aw-cuadrado-menos-identidad-pi}
\end{equation}
Products between distinct rows and columns vanish modulo three,
whereas each diagonal product equals two. One application effects a ternary
quarter-turn; two applications effect the oriented inversion.

\subsection{The slope 1/5 forced by the \(5\) regions}

Let
\[
 \lambda=(\lambda_1,\ldots,\lambda_5)^{\mathsf T}
 \in\mathbb Q_{\geq0}^5,
 \qquad
 \mathbf1^{\mathsf T}\lambda=1.
\]
Forgetting the regional label applies the symmetrizer
\[
 P_5=\frac15\mathbf1\mathbf1^{\mathsf T}.
\]
For every normalized vector,
\[
 P_5\lambda=\frac15\mathbf1.
\]
The invariance \(P_5\lambda=\lambda\) therefore has a unique solution:
\begin{equation}
 \boxed{\lambda=\frac15\mathbf1,\qquad t_1=\frac15.}
 \label{p12-83-c9-s1:eq:pendiente-primitiva-pi}
\end{equation}
The integer five comes from the regional microfiber; it is not recognized in a
formula for \(\pi\).

\subsection{Rational composition and compensator \texorpdfstring{\(239\)}{239}}

The law of slope composition is
\[
 t\star s:=\frac{t+s}{1-ts}.
\]
The first duplication gives
\[
 t_2=t_1\star t_1
 =\frac{2/5}{1-1/25}
 =\frac5{12}.
\]
The second gives
\[
 t_4=t_2\star t_2
 =\frac{10/12}{1-25/144}
 =\frac{120}{119}.
\]
The return to the unit diagonal requires \(q>0\) such that
\[
 \frac{120/119-1/q}{1+(120/119)/q}=1.
\]
Multiplying by \(119q\) produces
\[
 120q-119=119q+120,
\]
and forces
\begin{equation}
 \boxed{q=239.}
 \label{p12-83-c9-s1:eq:compensador-239-pi}
\end{equation}

\subsection{Gaussian identity and 1st half-turn}

We define
\begin{equation}
 \Theta_\pi
 :=16\arctan\frac15-4\arctan\frac1{239}.
 \label{p12-83-c9-s1:eq:theta-pi-autonoma}
\end{equation}
The successive squares are
\[
\begin{array}{c|r@{}l}
\text{power}&\multicolumn{2}{c}{\text{Gaussian integer}}\\ \hline
(5+i)^2&24&+10i\\
(5+i)^4&476&+480i\\
(5+i)^8&-3824&+456960i\\
(5+i)^{16}&-208797818624&-3494830080i\\
(239-i)^2&57120&-478i\\
(239-i)^4&3262465916&-54606720i
\end{array}
\]
and the final multiplication gives
\begin{equation}
 \boxed{
 (5+i)^{16}(239-i)^4
 =-681386607803576157184.}
 \label{p12-83-c9-s1:eq:identidad-gaussiana-pi-autonoma}
\end{equation}
The imaginary part is zero and the real part is negative. Therefore,
\[
 \Theta_\pi\equiv\pi\pmod{2\pi}.
\]
Para \(0<x<1\),
\[
 x-\frac{x^3}{3}<\arctan x<x.
\]
From here
\[
 \Theta_\pi>
 16\left(\frac15-\frac1{3\cdot5^3}\right)-\frac4{239}>3
\]
and
\[
 \Theta_\pi<\frac{16}{5}<4.
\]
The only negative orientation in that interval is the first positive half-turn:
\begin{equation}
 \boxed{\Theta_\pi=\pi.}
 \label{p12-83-c9-s1:eq:reconocimiento-pi-exacto}
\end{equation}

The order of subsequent Archimedean recognition is
\[
 \text{five regions}
 \to\frac15
 \to\frac5{12}
 \to\frac{120}{119}
 \to239
 \to\Theta_\pi
 \to\pi.
\]

\subsection{Arbitrary precision rational horocycles}

For \(0<x<1\), let
\[
 S_n(x)=\sum_{r=0}^{n}(-1)^r\frac{x^{2r+1}}{2r+1}.
\]
Leibniz's criterion gives
\[
 S_{2m+1}(x)<\arctan x<S_{2m}(x).
\]
We define
\[
 a_m^-=S_{2m+1}(1/5),\quad a_m^+=S_{2m}(1/5),
\]
\[
 b_m^-=S_{2m+1}(1/239),\quad b_m^+=S_{2m}(1/239),
\]
and
\begin{equation}
 \ell_{\pi,m}=16a_m^- -4b_m^+,
 \qquad
 h_{\pi,m}=16a_m^+ -4b_m^-.
 \label{p12-83-c9-s1:eq:horquilla-pi-autonoma}
\end{equation}
Then
\[
 \ell_{\pi,m}<\pi<h_{\pi,m}
\]
and
\begin{equation}
 h_{\pi,m}-\ell_{\pi,m}
 =
 \frac{16}{(4m+3)5^{4m+3}}
 +\frac{4}{(4m+3)239^{4m+3}}
 \longrightarrow0.
 \label{p12-83-c9-s1:eq:anchura-horquilla-pi}
\end{equation}
The lower bounds grow and the upper bounds decrease.

Since \(3<\pi<4\), the ternary reader acts on
\[
 r_\pi:=\pi-3\in(0,1),
\]
with bracket
\[
 [\ell_{\pi,m}-3,h_{\pi,m}-3].
\]

\subsection{Subsequent Archimedean certification of the already generated fiber}

Suppose the internal relation \(\operatorname{Ext}\) and coinductive
survival have already produced, without using Machin or a tabulated expansion,
the compatible extension \(N_{K+1}\) of the prefix \(N_K\), of length
\(L_K=6K\). The \(729\) subcylinders of its ambient fiber are
\[
 I_{K,u}=
 \left[
 \frac{729N_K+u}{3^{L_K+6}},
 \frac{729N_K+u+1}{3^{L_K+6}}
 \right),
 \qquad0\leq u<729.
\]
The order \(m=m(K)\) is then increased until the Leibniz bracket
locates the already produced subcylinder. Let
\[
 \ell_{\pi,K}:=\ell_{\pi,m(K)}-3,
 \qquad
 h_{\pi,K}:=h_{\pi,m(K)}-3.
\]
The compatible integer endpoints are
\[
 j_{\min}=
 \max\!\left(
 729N_K,\left\lfloor
 \ell_{\pi,K}3^{L_K+6}\right\rfloor
 \right),
\]
\[
 j_{\max}=
 \min\!\left(
 729N_K+728,\left\lceil
 h_{\pi,K}3^{L_K+6}\right\rceil-1
 \right).
\]
The condition of unitary localization is
\begin{equation}
 \boxed{j_{\min}=j_{\max},\qquad N_{K+1}=j_{\min}.}
 \label{p12-83-c9-s1:eq:seleccion-unitaria-subcilindro-pi}
\end{equation}
The first equality certifies unit localization; the second verifies
that this localization agrees with the internally generated extension. Neither
defines the TPK transition or selects among the \(729\) states.
Existence and uniqueness come from \(\operatorname{Ext}\) and
survival; Machin and the Leibniz brackets are subsequent
recognitions.

\subsection{Coinductive prolongation of the section of \(\pi\) through the inverse system of surviving histories}

Let \(\mathcal H_m^\pi\) be the set of enriched histories at depth
\(m\), and let
\[
 p_{n,m}:\mathcal H_n^\pi\longrightarrow\mathcal H_m^\pi
\]
the truncation. We define
\[
 \operatorname{Surv}_m^{(N),\pi}
 :=p_{N,m}(\mathcal H_N^\pi),
\]
\[
 \operatorname{Surv}_m^\pi
 :=\bigcap_{N\geq m}\operatorname{Surv}_m^{(N),\pi}.
\]
When the relation \(\operatorname{Ext}\) is nonempty, single-valued and
compatible at every level, the sets are singletons and satisfy
\[
 p_{m+1,m}(\operatorname{Surv}_{m+1}^{\pi})
 =\operatorname{Surv}_{m}^{\pi}.
\]
In that domain of extendability, there exists, therefore, a unique section.
\[
 x_\pi\in
 X_\pi:=\varprojlim_m\operatorname{Surv}_m^\pi.
\]
The cylinders are nested and their diameters are \(3^{-6K}\), which tends to
zero. Their intersection contains a unique point, identified by the angular
character with \(r_\pi=\pi-3\).

For every decimal depth \(M\), one may choose a level \(K(M)\) whose
cylinder lies within a single decimal cell of diameter \(10^{-M}\).
The deep prefixes truncate exactly to the preceding ones.

\subsection{Characterization of the coinductive channel of \(\pi\): selection of \(w_6^\pi\), \(5\) regions, oriented inversion, and uniqueness of the limit section}

\begin{theorem}[Characterization of the closure channel]
\label{p12-83-c9-s1:thm:caracterizacion-final-pi-autonoma}
Starting from the catalog, the action \(D_3\), the oriented gauge, the
operators \(L_0,L_1,A_W\), the slope law and the rational filtration:
\begin{enumerate}
 \item \(w_6^\pi=010211\) is selected;
 \item five regions are constructed with the identities of
       \eqref{p12-83-c9-s1:eq:descomposiciones-cinco-regiones-pi};
 \item each region belongs to a minimum closure of length eight;
 \item \(A_W^2=-I_6\) fixes the oriented inversion;
 \item regional symmetrization forces \(1/5\);
 \item the composition forces \(120/119\);
 \item the unit return forces \(239\);
 \item the Gaussian identity recognizes the first half-turn;
 \item the brackets and the fiber of \(729\) elements produce a section
       compatible at every finite depth.
\end{enumerate}
Its evaluation is
\[
 \boxed{\pi=16\arctan\frac15-4\arctan\frac1{239}.}
\]
\end{theorem}

The complete chain is
\[
 \mathrm{APP}\to\mathrm{TRIT}\to\mathrm{TPK}
 \to104\,976\to468\to243\to w_6^\pi
 \to w_{30}^\pi\to R_{36}\to\Gamma_9
 \to X_\pi\to\pi.
\]

\begin{criterion}[Channel Closure Falsification]
The derivation is refuted if:
\begin{enumerate}
 \item the closure orbit ceases to be unique;
 \item one of the five regions disappears or its complete identity changes;
 \item one of the five closures has length less than eight;
 \item \(A_W^2\ne2I_6\);
 \item the symmetrizer admits another normalized distribution;
 \item the compositions do not produce \(5/12\), \(120/119\), and \(239\);
 \item the Gaussian identity fails;
 \item a bracket ceases to contain the angular character;
 \item a level does not traverse the \(729\) elements of the ambient fiber or does not
       locate exactly one extension;
 \item a target expansion intervenes before recognition.
\end{enumerate}
\end{criterion}
\fi

The complete construction of the closure character, the five regions, and the
structural biography of \(\pi\) was established in
\cref{chap:83-08}. Starting from that same already constructed character, the present
chapter now prolongs the refinement toward \(e\) and \(\varphi\), without
restarting the genealogy or repeating its first stage.


% Unidad U013. Biografía estructural completa del carácter exponencial.

\section{Propagation, memory return and construction of the exponential character}
\label{p12-83-c9-s2:chap:biografia-completa-e}

\subsection{Finite orbital selector: existence and uniqueness of compatible emission}

Let \(\mathcal W=\mathbb F_3^6\). For an emission
\(U=(u_1,\ldots,u_6)\), the visible word is
\[
 \Psi(U)=((-u_1)\bmod3,\ldots,(-u_6)\bmod3).
\]
On \(\mathcal W\), the dihedral action is given by
\[
\begin{aligned}
 r(u_1,u_2,u_3,u_4,u_5,u_6)
 &=(u_3,u_1,u_2,u_5,u_6,u_4),\\
 s(u_1,u_2,u_3,u_4,u_5,u_6)
 &=(u_4,u_5,u_6,u_1,u_2,u_3).
\end{aligned}
\]
For \(w\in\mathcal W\), define
\[
 n(w)=\#\Psi^{-1}(w),
 \qquad
 M(w)=\sum_{U\in\Psi^{-1}(w)}\operatorname{mult}(U),
\]
and
\[
 \nu(w)=
 \bigl(
 \#\{j:w_j=0\},
 \#\{j:w_j=1\},
 \#\{j:w_j=2\}
 \bigr).
\]

\begin{theorem}[Structural selection of the exponential channel]
\label{p12-83-c9-s2:thm:seleccion-estructural-e}
Among the \(43\) orbits of the visible catalog, there is a unique orbit whose
fibers simultaneously satisfy:
\[
 n(w)=1,\qquad M(w)=144,\qquad
 \mathfrak D(w)=\{0,3,6\},\qquad
 \nu(w)=(2,3,1).
\]
The gauge
\[
 \operatorname{diag}_c=6,
 \qquad \operatorname{card}=ES
\]
select in it
\[
 \boxed{w_6^e=201101.}
\]
\end{theorem}

\begin{proof}
The four conditions are evaluated over the \(243\) words in the catalog, not over a sample. A single orbit satisfies the joint predicate. In that orbit, exactly one emission has both diagonal \(6\) and orientation \(ES\). The enumeration uses only fields from the APP--TRIT catalog and does not consult an expansion of \(e\).
\end{proof}

\subsection{Visible lift and chain of blocks}

The recurrence \(L_0,L_0,L_1,L_1\) produces
\[
 w_{30}^{e}
 =201101|121221|102011|012222|102011,
\]
and the joint boundary contributes
\[
 u_5^e=021222.
\]
The record of thirty-six trits remains
\[
 w_{36}^{e}
 =201101|121221|102011|012222|102011|021222.
\]

\subsection{Decimal signature and positional square}

The initial reading of twelve digits is
\[
 \mathbf d_e^{(12)}
 =718281828459
 =7\,\underbrace{1828\,1828}_{\text{local square}}\,459.
\]
The central equality is the concatenation
\[
 1828\mathbin{\|}1828,
\]
not an arithmetic product nor the supposed beginning of a period.

\begin{definition}[Positional square]
A word \(d_1\cdots d_{12}\) contains a square of length \(\ell\)
at position \(p\) when
\[
 d_p\cdots d_{p+\ell-1}
 =d_{p+\ell}\cdots d_{p+2\ell-1}.
\]
The position and the lateral contexts belong to the signature.
\end{definition}

\begin{table}[htbp]
\centering
\caption{Comprehensive census of squares in twelve-digit readings.}
\label{p12-83-c9-s2:tab:censo-cuadrados-e}
\begin{tabular}{c*{6}{r}}
\toprule
Length \(\ell\)&1&2&3&4&5&6\\
\midrule
Occurrences&240&21&3&2&0&0\\
Words that contain them&153&19&3&2&0&0\\
\bottomrule
\end{tabular}
\end{table}

The other word with a square of length four is
\[
 575157518472
 =\underbrace{5751\,5751}_{\text{initial position }1}\,8472.
\]

\begin{proposition}[Positional uniqueness]
\label{p12-83-c9-s2:prop:cuadrado-posicional-e}
The reading of channel \(e\) is the only one that has a square of length
four after a prefix of length one and with contexts \(7\) and \(459\).
\end{proposition}

\begin{proof}
The census leaves only two words with squares of length four. One starts
it at the first position; the other after the prefix \(7\). The
position and contexts separate the two possibilities.
\end{proof}

\subsection{Visible return and memory rupture}

In the elevation
\[
 201101|121221|102011|012222|102011
\]
the third and the fifth blocks coincide:
\[
 B_3=B_5=102011.
\]
Their prestates modulo \(27\) are, however,
\[
 p_3=(25,11,7,17,8,16),
 \qquad
 p_5=(7,8,4,23,26,7),
\]
and the subsequent quotients:
\[
 q_3=(6,13,9,2,13,1),
 \qquad
 q_5=(15,0,3,19,23,25).
\]
\(p_3\ne p_5\) and \(q_3\ne q_5\) hold.

\begin{figure}[htbp]
\centering
\begin{tikzpicture}[
 node distance=11mm and 7mm,
 every node/.style={font=\small},
 block/.style={draw,rounded corners,minimum width=17mm,minimum height=8mm},
 state/.style={draw, dashed, rounded corners, align=center, inner sep=3pt},
 >=stealth
]
\node[block] (b1) {\(201101\)};
\node[block,right=of b1] (b2) {\(121221\)};
\node[block,right=of b2] (b3) {\(102011\)};
\node[block,right=of b3] (b4) {\(012222\)};
\node[block,right=of b4] (b5) {\(102011\)};
\draw[->] (b1)--(b2);
\draw[->] (b2)--(b3);
\draw[->] (b3)--(b4);
\draw[->] (b4)--(b5);
\node[state,below=of b3] (p3)
 {\(p_3=(25,11,7,17,8,16)\)\\
  \(q_3=(6,13,9,2,13,1)\)};
\node[state,below=of b5] (p5)
 {\(p_5=(7,8,4,23,26,7)\)\\
  \(q_5=(15,0,3,19,23,25)\)};
\draw[->] (p3)--(b3);
\draw[->] (p5)--(b5);
\draw[<->,thick] (b3) to[bend left=28]
 node[above]{same word visible} (b5);
\end{tikzpicture}
\caption{Visible return of \(102011\) with renewal of the prestate and the
quotient.}
\label{p12-83-c9-s2:fig:ruptura-memoria-e}
\end{figure}

\begin{theorem}[Visible-phase return without periodicity of the state]
\label{p12-83-c9-s2:thm:retorno-visible-no-periodico-e}
\(B_3=B_5\) is a return of the visible projection, not a periodic orbit
of the enriched state.
\end{theorem}

\begin{proof}
The return of the complete state would require equality of all its coordinates.
The prestate and quotient inequalities exclude it. Nor does
\(1828|1828\) begin a decimal period: after the second block,
\(4\) appears, whereas a periodic continuation would require \(1\).
\end{proof}

\subsection{Subsequent factorial certification of the coinductive exponential section}

Let \(\mathcal A\) be a unital commutative algebra over \(\mathbb Q\), and
let \(P_t\in\mathcal A[[t]]\) be the formal propagator. Temporal
concatenation imposes
\[
 P_{s+t}=P_sP_t,
 \qquad
 P_0=1.
\]
We write
\[
 P_t=\sum_{n=0}^{\infty}a_nt^n.
\]
The unit fixes \(a_0=1\), and normalized elementary transport fixes
\(a_1=1\).

\begin{theorem}[Factorial recurrence]
\label{p12-83-c9-s2:thm:recurrencia-factorial-e}
\[
 (n+1)a_{n+1}=a_n
 \qquad(n\geq0),
\]
and, therefore,
\[
 a_n=\frac1{n!},
 \qquad
 P_t=\sum_{n=0}^{\infty}\frac{t^n}{n!}.
\]
En \(t=1\),
\[
 \boxed{P_1=e.}
\]
\end{theorem}

\begin{proof}
The coefficient of \(s^nt\) in \(P_{s+t}\) is
\[
 \binom{n+1}{n}a_{n+1}=(n+1)a_{n+1}.
\]
In \(P_sP_t\) it is \(a_na_1=a_n\). Equality of series produces the
recurrence. Induction from \(a_0=1\) gives \(a_n=1/n!\).
\end{proof}

\begin{corollary}[Differential equation]
\[
 \frac{d}{dt}P_t=P_t,
 \qquad P_0=1.
\]
\end{corollary}

\begin{proof}
\[
 P_t'=\sum_{n\geq0}(n+1)a_{n+1}t^n
 =\sum_{n\geq0}a_nt^n=P_t.
\]
\end{proof}

\subsection{Cofinal Forking of Rational Tangles and Convergence to the Published Value}

Let
\[
 S_N=\sum_{n=0}^{N}\frac1{n!}.
\]

\begin{theorem}[Explicit factorial bound]
\label{p12-83-c9-s2:thm:cota-factorial-e}
For \(N\geq1\),
\begin{equation}
 \boxed{
 S_N<e<
 S_N+\frac{N+2}{(N+1)(N+1)!}.}
 \label{p12-83-c9-s2:eq:horquilla-factorial-e}
\end{equation}
\end{theorem}

\begin{proof}
The rest is
\[
 R_N=\sum_{k=0}^{\infty}\frac1{(N+1+k)!}>0.
\]
Para \(k\geq0\),
\[
 (N+1+k)!
 \geq(N+1)!(N+2)^k.
\]
Therefore,
\[
 R_N
 \leq\frac1{(N+1)!}
 \sum_{k=0}^{\infty}\frac1{(N+2)^k}
 =\frac{N+2}{(N+1)(N+1)!}.
\]
The inequality is strict from \(k=2\) onward.
\end{proof}

We define
\[
 J_N^{(e)}=
 \left[
 S_N,\,
 S_N+\frac{N+2}{(N+1)(N+1)!}
 \right]
\]
and, for the fractional part,
\[
 \widetilde J_N^{(e)}=J_N^{(e)}-2.
\]
The width converges to zero. For each \(M\), there exists \(N\) such that
\[
 \operatorname{diam}J_N^{(e)}<10^{-M}.
\]
If the endpoints belong to the same decimal cell, the first \(M\)
digits are certified by rational arithmetic.

\subsection{Coinductive prolongation of the exponential section through ternary subcylinders of the Topological Phase Kernel}

At level \(K\), TPK and the relation \(\operatorname{Ext}\) have already generated
a unique compatible subcylinder. The smallest \(N=N_e(K)\) for which
the factorial bracket is contained in that subcylinder is taken. \(N_e(K)\) is an
order of analytic certification, not a generating index or a rule for
selecting the fiber.

The TPK cylinders are compatible under truncation independently of the
semigroup. The semigroup theorem subsequently identifies their Archimedean
evaluation with \(e-2\); the integer part recovers \(e\). If the orbit,
\(\operatorname{Ext}\) or truncation fails, generation fails; if the
recurrence or the factorial bound fails, this analytic certificate fails, not the
APP--TRIT--TPK genealogy.

\subsection{Adversarial certificate for the generation of \(e\): orbital uniqueness, factorial recurrence, rational brackets and cylindrical prolongation}

\begin{criterion}[Exponential Channel Falsification]
The construction is refuted if:
\begin{enumerate}
 \item another orbit satisfies the complete finite predicate;
 \item the signature \(7[1828\,1828]459\) is not positionally unique;
 \item the prestates or quotients of the two blocks \(102011\) coincide;
 \item the semigroup does not force \((n+1)a_{n+1}=a_n\);
 \item the bound \eqref{p12-83-c9-s2:eq:horquilla-factorial-e} fails;
 \item a level does not locate a compatible extension uniquely;
 \item a tabulated expansion intervenes before recognition.
\end{enumerate}
\end{criterion}


% Unidad U014. Biografía estructural completa del carácter de autoescala.

\section{Modal incidence, Perron ray, and construction of the self-similar character}
\label{p12-83-c9-s3:chap:biografia-completa-phi}

\subsection{Finite Selection of the Self-Scaling Seed}

Let \(\mathcal U_6\subset\{0,\ldots,999\}^6\) be the catalog of the
\(468\) distinct emissions constructed in
\cref{chap:semillas-emisor-54}, and let
\[
 \Psi:\mathcal U_6\longrightarrow\mathbb F_3^6,
 \qquad
 \Psi(U)=((-U_j)\bmod3)_{j=1}^{6}.
\]
For \(w\in\Psi(\mathcal U_6)\), we recall the invariants
\[
 n(w)=\#\Psi^{-1}(w),
 \qquad
 M(w)=\sum_{U\in\Psi^{-1}(w)}\operatorname{mult}(U),
\]
and the occupation vector
\[
 \nu(w)=
 \bigl(\#\{j:w_j=0\},\#\{j:w_j=1\},\#\{j:w_j=2\}\bigr).
\]

\begin{theorem}[Orbital selection of the self-scale]
\label{p12-83-c9-s3:thm:seleccion-orbital-phi}
The dihedral action on \(\Psi(\mathcal U_6)\) contains a unique orbit
that satisfiis the self-scale predicate
\[
 n(w)=1,
 \qquad
 M(w)=144,
 \qquad
 \nu(w)=(2,2,2).
\]
The internal caliber
\[
 \operatorname{diag}_c=6,
 \qquad
 \operatorname{card}=ES
\]
select in that orbit the word
\[
 \boxed{w_6^{\varphi}=121200.}
\]
Its only emission is
\[
 \boxed{U_{\varphi}=(830,943,830,109,252,252),}
\]
of multiplicity \(144\), multiplicative signature \(555933\), additive
orientation \(ES\), multiplicative support of cardinality six and residual type
\(A\).
\end{theorem}

\begin{proof}
The direct projection gives
\[
 U_{\varphi}\bmod3=(2,1,2,1,0,0),
 \qquad
 -U_{\varphi}\bmod3=(1,2,1,2,0,0).
\]
The enumeration is carried out over the \(243\) visible words and their
\(43\) orbits, not over a selection of examples. The joint predicate
\(n=1\), \(M=144\), \(\nu=(2,2,2)\) leaves one orbit. Within it, the
declared diagonal and orientation leave one emission. At this stage neither
the equation \(x^2-x-1=0\), nor a decimal table, nor the
conventional name of the constant has been introduced.
\end{proof}

\subsection{Blockwise lift and 1st coinductive boundary}

Let \(b_1=w_6^{\varphi}\), and let
\(\varepsilon=(0,0,1,1)\). With the matrices \(L_0,L_1\) constructed in
\cref{hmtctx:p12:chap:orbitas-elevacion-r36}, we define
\[
 b_{k+1}=b_kL_{\varepsilon_k}\pmod3
 \qquad(1\leq k\leq4),
\]
and concatenate
\[
 w_{6(k+1)}^{\varphi}=w_{6k}^{\varphi}\mathbin{\|}b_{k+1}.
\]
Matrix calculus produces, block by block,
\[
\begin{aligned}
 b_1&=121200,\\
 b_2&=112202,\\
 b_3&=121020,\\
 b_4&=010210,\\
 b_5&=010200,
\end{aligned}
\]
and, therefore,
\begin{equation}
 \boxed{
 w_{30}^{\varphi}
 =121200|112202|121020|010210|010200.}
 \label{p12-83-c9-s3:eq:w30-phi}
\end{equation}
The dual neutrality of the boundary adds
\[
 u_5^{\varphi}=102011,
\]
so that
\begin{equation}
 \boxed{
 w_{36}^{\varphi}
 =121200|112202|121020|010210|010200|102011.}
 \label{p12-83-c9-s3:eq:w36-phi}
\end{equation}

\begin{remark}[Two distinct uses of a visible word]
The word \(102011\) is also visible in the biography of \(e\). The
equality of a projection does not identify enriched states: the channel,
prefix, quotient, carry, phase, and memory belong to the type of the
object. This separation is mandatory before any
real character is constructed.
\end{remark}

\subsection{From the tritic calendar to the incidence multigraph}

\subsubsection{Primitive modal incidence \(F_{\rm av}\) and memory dynamics of order \(1\)}

The primitive triticale calendar is
\[
 (+1,-1,0).
\]
We fix two visible modis, \(\Sigma\) and \(\Pi\), and the rule
\begin{equation}
 +1:m\longmapsto\Sigma,
 \qquad
 -1:m\longmapsto\Pi,
 \qquad
 0:m\longmapsto m.
 \label{p12-83-c9-s3:eq:regla-modal-primitiva-phi}
\end{equation}
Under the cyclic closure condition, the modal orbit produced by
\eqref{p12-83-c9-s3:eq:regla-modal-primitiva-phi} is
\[
 \boxed{(\Sigma,\Pi,\Pi).}
\]
The three consecutive pairs are
\[
 \Sigma\longrightarrow\Pi,
 \qquad
 \Pi\longrightarrow\Pi,
 \qquad
 \Pi\longrightarrow\Sigma.
\]

\begin{definition}[Chronological incidence matrix]
For \(i,j\in\{\Sigma,\Pi\}\), let
\[
 N_3(i,j)=
 \#\{r\in\mathbb Z/3\mathbb Z:(m_r,m_{r+1})=(i,j)\}.
\]
The subscript \(3\) counts instants of the primitive calendar; it does not denote a
matrix power.
\end{definition}

\begin{theorem}[Necessary descent of modal incidence]
\label{p12-83-c9-s3:thm:descenso-incidencia-modal-phi}
In the order \((\Sigma,\Pi)\),
\begin{equation}
 \boxed{
 N_3=
 \begin{pmatrix}
 0&1\\
 1&1
 \end{pmatrix}.}
 \label{p12-83-c9-s3:eq:N3-phi}
\end{equation}
By exact repetition of the calendar,
\begin{equation}
 \boxed{
 N_9=3N_3,
 \qquad
 N_{27}=9N_3,
 \qquad
 N_{108}=36N_3.}
 \label{p12-83-c9-s3:eq:incidencias-9-27-108-phi}
\end{equation}
These three equalities are equalities of edge counts; they do not assert
\(N_9=N_3^3\), \(N_{27}=N_3^9\) or \(N_{108}=N_3^{36}\).
\end{theorem}

\begin{proof}
\(\Sigma\to\Sigma\) does not occur; each of the other three pairs occurs
once. This determines the four entries of \(N_3\). The calendars of
lengths \(9\), \(27\), and \(108\) contain respectively \(3\), \(9\),
and \(36\) copies of the primitive block. Each repetition multiplies all the
counts by the same integer, which proves
\eqref{p12-83-c9-s3:eq:incidencias-9-27-108-phi}.
\end{proof}

\subsubsection{Publication areal–volumetrica and razon estable \(\varphi^{-2}\) under the measure of Parry}

The modal label and the geometric sheet are distinct objects. The publication chart is
\[
 \mathfrak g_{\rm av}(\Sigma)=\mathscr H_{\rm ar},
 \qquad
 \mathfrak g_{\rm av}(\Pi)=\mathscr H_{\rm vol}.
\]
Transporting \eqref{p12-83-c9-s3:eq:N3-phi} through \(\mathfrak g_{\rm av}\), we obtain
the incidence operator
\begin{equation}
 \boxed{
 F_{\rm av}=
 \begin{pmatrix}
 0&1\\
 1&1
 \end{pmatrix}.}
 \label{p12-83-c9-s3:eq:Fav-derivada-phi}
\end{equation}
Thus, the two cross-transitions, the unique volumetric self-retention, and the
absence of areal self-retention are derived from the calendar; they are not four
independent inputs.

\begin{remark}[Scope of the descent]
Forgetting the elementary phase does not by itself produce a strong Markov
quotient of the Topological Phase Kernel, because the visible mode does not determine
which of the three internal operators will act at the next step. What
does descend without an additional hypothesis is the edge multigraph of a
complete block. Its minimal memory-one envelope has adjacency
matrix \(F_{\rm av}\).
\end{remark}

\subsection{Invariant positive ray and uniqueness of self-scaling}

\begin{theorem}[Perron ray derived]
\label{p12-83-c9-s3:thm:rayo-perron-phi}
The positive cone of \(F_{\rm av}\) contains a unique eigenray. If we
normalize it as \(v=(1,x)^{\mathsf T}\), then
\[
 x^2-x-1=0,
 \qquad
 x>0,
\]
and, therefore,
\begin{equation}
 \boxed{x=\frac{1+\sqrt5}{2}=\varphi.}
 \label{p12-83-c9-s3:eq:phi-perron}
\end{equation}
\end{theorem}

\begin{proof}
The equation \(F_{\rm av}v=\lambda v\) gives, coordinate by coordinate,
\[
 x=\lambda,
 \qquad
 1+x=\lambda x.
\]
Eliminating \(\lambda\) yields \(x^2-x-1=0\). Its roots are \((1\pm\sqrt5)/2\), and only the positive root lies in the interior of the cone. The matrix \(F_{\rm av}\) is primitive, since \(F_{\rm av}^2\) has all its entries positive; the Perron--Frobenius theorem ensures that this positive ray is simple and unique. The name \(\varphi\) is assigned only after this construction.
\end{proof}

\begin{corollary}[Scalar Self-Similarity Equation]
\[
 \varphi=1+\frac1\varphi,
 \qquad
 \varphi^2=\varphi+1,
 \qquad
 \varphi^{-1}=\varphi-1.
\]
Thi identitiis are consequencis of \eqref{p12-83-c9-s3:eq:phi-perron}; they select neither
the seed nor the matrix.
\end{corollary}

\subsection{Fibonacci Recurrence of the Perron Ray and Diophantine Embedding of Rational Gaps}

Let \(F_0=0\), \(F_1=1\), and
\[
 F_{n+1}=F_n+F_{n-1}\qquad(n\geq1).
\]

\begin{theorem}[Exact Powers of Incidence]
\label{p12-83-c9-s3:thm:potencias-fibonacci-phi}
For every \(n\geq1\),
\begin{equation}
 \boxed{
 F_{\rm av}^{\,n}=
 \begin{pmatrix}
 F_{n-1}&F_n\\
 F_n&F_{n+1}
 \end{pmatrix}.}
 \label{p12-83-c9-s3:eq:potencias-Fav-fibonacci}
\end{equation}
\end{theorem}

\begin{proof}
For \(n=1\), the identity is the definition of \(F_{\rm av}\). If it holds for \(n\), right multiplication by \(F_{\rm av}\) gives
\[
 \begin{pmatrix}
 F_n&F_{n-1}+F_n\\
 F_{n+1}&F_n+F_{n+1}
 \end{pmatrix}
 =
 \begin{pmatrix}
 F_n&F_{n+1}\\
 F_{n+1}&F_{n+2}
 \end{pmatrix},
\]
which is the case \(n+1\).
\end{proof}

\begin{theorem}[Cassini's Identity]
\label{p12-83-c9-s3:thm:cassini-phi}
For every \(n\geq1\),
\begin{equation}
 \boxed{F_{n+1}F_{n-1}-F_n^2=(-1)^n.}
 \label{p12-83-c9-s3:eq:cassini-phi}
\end{equation}
\end{theorem}

\begin{proof}
The determinant of \(F_{\rm av}\) is \(-1\). Taking determinants in
\eqref{p12-83-c9-s3:eq:potencias-Fav-fibonacci},
\[
 F_{n-1}F_{n+1}-F_n^2
 =\det(F_{\rm av}^{\,n})=(-1)^n.
\]
\end{proof}

We define \(r_n=F_{n+1}/F_n\) for \(n\geq1\). From Cassini it follows that
\[
 r_n-r_{n-1}
 =\frac{F_{n+1}F_{n-1}-F_n^2}{F_nF_{n-1}}
 =\frac{(-1)^n}{F_nF_{n-1}}.
\]
Therefore, the approximations alternate around the Perron ray.

\begin{proposition}[Perron Rational Gaps]
\label{p12-83-c9-s3:prop:horquillas-perron-phi}
For \(m\geq1\),
\begin{equation}
 \frac{F_{2m+2}}{F_{2m+1}}
 <\varphi<
 \frac{F_{2m+1}}{F_{2m}},
 \label{p12-83-c9-s3:eq:horquilla-fibonacci-phi}
\end{equation}
and the width is
\begin{equation}
 \boxed{
 \frac{F_{2m+1}}{F_{2m}}
 -\frac{F_{2m+2}}{F_{2m+1}}
 =\frac{1}{F_{2m}F_{2m+1}}.}
 \label{p12-83-c9-s3:eq:anchura-fibonacci-phi}
\end{equation}
\end{proposition}

\begin{proof}
The preceding alternation places the even quotients below and the odd ones
above \(\varphi\). The formula for the width is
\eqref{p12-83-c9-s3:eq:cassini-phi} with \(n=2m+1\), after putting both fractions over a
common denominator. Since \(F_n\to\infty\), the width tends to zero.
\end{proof}

\subsection{Memory incidence of order \(1\) and Parry measure of the symbolic envelope}

Let \(X_{\rm av}\) be the smallest subshift of finite type over
\(\{\mathscr H_{\rm ar},\mathscr H_{\rm vol}\}\) that admits exactly
the three pairs observed in the modal calendar. Its adjacency matrix is
\(F_{\rm av}\).

\begin{theorem}[Transition and Parry Measure]
\label{p12-83-c9-s3:thm:parry-phi}
Let \(r=(1,\varphi)^{\mathsf T}\). The Parry stochastic matrix is
\begin{equation}
 \boxed{
 P_{ij}=
 \frac{(F_{\rm av})_{ij}r_j}{\varphi r_i}
 =
 \begin{pmatrix}
 0&1\\
 \varphi^{-2}&\varphi^{-1}
 \end{pmatrix}.}
 \label{p12-83-c9-s3:eq:parry-matriz-phi}
\end{equation}
The measure estacionaria of vertices is proporcional to
\((1,\varphi^2)\), hence
\begin{equation}
 \boxed{
 \frac{\mu_{\rm ar}}{\mu_{\rm vol}}=\varphi^{-2}.}
 \label{p12-83-c9-s3:eq:parry-razon-phi}
\end{equation}
\end{theorem}

\begin{proof}
The sum of row \(i\) of \(P\) is
\[
 \frac{(F_{\rm av}r)_i}{\varphi r_i}=1.
\]
Since \(F_{\rm av}\) is symmetric, the left and right eigenvectors
coincide. The stationary weights are proportional to their coordinatewise
products, that is, to \((1,\varphi^2)\). The ratio between the
two weights is \(\varphi^{-2}\).
\end{proof}

\begin{remark}[Three measures that must not be confused]
The chronological frequency of the primitive orbit is
\[
 \nu_{\rm cyc}(\mathscr H_{\rm ar})=\frac13,
 \qquad
 \nu_{\rm cyc}(\mathscr H_{\rm vol})=\frac23.
\]
The uniform measure \(\tau_{468}\) lives on the emission catalog. The
Parry measure lives on admissible histories of arbitrary length. The
first counts instants, the second counts emissions, and the third weights
dynamic cylinders. None is an approximation of the others.
\end{remark}

\subsection{Multiplicative character and cylindrical conservation}

For an admissible route \(\gamma:i\to j\) of length \(n\), we define
\begin{equation}
 \chi_P(\gamma)=\varphi^n\frac{r_i}{r_j}.
 \label{p12-83-c9-s3:eq:caracter-parry-phi}
\end{equation}

\begin{proposition}[Concatenation multiplicativity]
\label{p12-83-c9-s3:prop:multiplicatividad-parry-phi}
If \(\gamma_1:i\to j\) and \(\gamma_2:j\to k\), then
\[
 \chi_P(\gamma_2\circ\gamma_1)
 =\chi_P(\gamma_2)\chi_P(\gamma_1).
\]
On to closed route of length \(n\),
\[
 \chi_P(\gamma)=\varphi^n.
\]
\end{proposition}

\begin{proof}
Length adds and the intermediate coordinate telescopes:
\[
 \varphi^{n_1+n_2}\frac{r_i}{r_k}
 =\left(\varphi^{n_1}\frac{r_i}{r_j}\right)
  \left(\varphi^{n_2}\frac{r_j}{r_k}\right).
\]
If \(i=k\), the quotient terminal is one.
\end{proof}

We normalize the left vector as
\[
 \ell=\frac{r}{1+\varphi^2},
 \qquad
 \ell^{\mathsf T}r=1.
\]
The measure of a cylinder is
\[
 \mu[x_0\cdots x_n]
 =\ell_{x_0}r_{x_n}\varphi^{-n}.
\]
Therefore,
\begin{equation}
 V_n(x)
 :=\frac{\mu[x_0]}{\mu[x_0\cdots x_n]}
 =\varphi^n\frac{r_{x_0}}{r_{x_n}}
 =\chi_P(x_0\cdots x_n).
 \label{p12-83-c9-s3:eq:volumen-cilindrico-phi}
\end{equation}
For to typed dimension \(D>0\), the scale
\[
 a_n^{(D)}(x)=V_n(x)^{1/D}
\]
satisfies the exact law
\begin{equation}
 \boxed{
 \bigl(a_n^{(D)}(x)\bigr)^D
 \mu[x_0\cdots x_n]=\mu[x_0].}
 \label{p12-83-c9-s3:eq:conservacion-cilindrica-phi}
\end{equation}

In three dimensions and over closed returns,
\[
 \frac{a_9}{a_0}=\varphi^3,
 \qquad
 \frac{a_{27}}{a_0}=\varphi^9,
 \qquad
 \frac{a_{108}}{a_0}=\varphi^{36}.
\]
If \(A_k=a_{9k}\), then
\begin{equation}
 A_{k+1}-2A_k+A_{k-1}
 =A_{k-1}(\varphi^3-1)^2>0.
 \label{p12-83-c9-s3:eq:convexidad-retornos-phi}
\end{equation}
This is convexity in return depth; it is not interpreted as a
kinematic quantity without an additional duration chart.

\subsection{Quadratic memory and contracted branch}

The eigenvalues of \(F_{\rm av}\) are
\[
 \varphi,
 \qquad
 -\varphi^{-1}.
\]
The contracted branch changes orientation at each iteration. To preserve
that sign before squaring, we pass to
\(\operatorname{Sym}^2(\mathbb R^2)\). In the basis
\((x^2,2xy,y^2)\),
\begin{equation}
 \operatorname{Sym}^2(F_{\rm av})=
 \begin{pmatrix}
 0&0&1\\
 0&1&2\\
 1&1&1
 \end{pmatrix}.
 \label{p12-83-c9-s3:eq:sym2-Fav-phi}
\end{equation}

\begin{theorem}[Spectrum of the second-order memory]
\label{p12-83-c9-s3:thm:espectro-sym2-phi}
\[
 \det\bigl(tI-\operatorname{Sym}^2(F_{\rm av})\bigr)
 =(t+1)(t^2-3t+1),
\]
and
\begin{equation}
 \boxed{
 \operatorname{Spec}\bigl(\operatorname{Sym}^2(F_{\rm av})\bigr)
 =\{\varphi^2,-1,\varphi^{-2}\}.}
 \label{p12-83-c9-s3:eq:espectro-sym2-phi}
\end{equation}
The only positive stable mode is \(\varphi^{-2}\).
\end{theorem}

\begin{proof}
The eigenvalues of a degree-two symmetric power are the products
\(\lambda_1^2\), \(\lambda_1\lambda_2\), and \(\lambda_2^2\). With
\(\lambda_1=\varphi\) and \(\lambda_2=-\varphi^{-1}\), one obtains
\(\varphi^2\), \(-1\), and \(\varphi^{-2}\). The first exceeds one, the
second has modulus one, and the third belongs to \((0,1)\).
\end{proof}

On a closed route with \(V_n=\varphi^n\), the stable mode satisfies
\begin{equation}
 \boxed{M_n=\varphi^{-2n}=V_n^{-2}.}
 \label{p12-83-c9-s3:eq:modo-estable-phi}
\end{equation}

\subsection{Normalized hyperbolic realization of the self-scale: \(\exp(2\log\varphi\,J_h)=F_{\mathrm{av}}^2\)}

Let
\[
 A=F_{\rm av}^{\,2}=
 \begin{pmatrix}1&1\\1&2\end{pmatrix},
 \qquad
 J_h=\frac{2A-3I}{\sqrt5}.
\]
Then
\[
 J_h^2=I,
 \qquad
 A=\frac32I+\frac{\sqrt5}{2}J_h.
\]
Since
\[
 \cosh(2\log\varphi)=\frac32,
 \qquad
 \sinh(2\log\varphi)=\frac{\sqrt5}{2},
\]
we obtain
\begin{equation}
 \boxed{
 \exp(2\log\varphi\,J_h)=F_{\rm av}^{\,2}.}
 \label{p12-83-c9-s3:eq:realizacion-hiperbolica-phi}
\end{equation}
The self-scaling constant is therefore simultaneously the Perron
eigenvalue of the incidence and the exponential parameter of its normalized
hyperbolic realization.

\subsection{Characterization, brackets and prolongation}

The finite construction determines the positive ray of \(F_{\rm av}\). The
subsequent analytic characterization is
\[
 x^2-x-1=0,
 \qquad x>0.
\]
The brackets \eqref{p12-83-c9-s3:eq:horquilla-fibonacci-phi} have rational endpoints
and explicit width \eqref{p12-83-c9-s3:eq:anchura-fibonacci-phi}. For each
nonadic level \(K\), we choo the least \(m=m_{\varphi}(K)\) who bracket remains
contained in a single subcylinder among the \(729\) elements of the
ambient fiber. Minimality converts localization into a function, not into to
tacit choice.

The selected cylinders are compatible with the truncations, and their
diameters tend to zero. Their intersection contains a unique real number; the
\cref{p12-83-c9-s3:thm:rayo-perron-phi} identifies it with \(\varphi\). The Fibonacci
quotients are therefore rational certificates of localization, not decimal data
used to decide the seed.

\subsection{Internal certificate of uniqueness, compatibility, and convergence of self-scaling}

\begin{criterion}[Refutation of the self-scale biography]
The construction remains refuted if any of the following facts
occurs:
\begin{enumerate}
 \item the finite predicate selects more than one orbit or more than one emission;
 \item the recurrence \(L_0,L_0,L_1,L_1\) does not produce
       \eqref{p12-83-c9-s3:eq:w30-phi};
 \item the calendar does not produce exactly the three pairs used in
       \eqref{p12-83-c9-s3:eq:N3-phi};
 \item a count multiple \(N_{9},N_{27},N_{108}\) is confused with a
       power of \(N_3\);
 \item \(F_{\rm av}\) has another positive eigenray;
 \item \eqref{p12-83-c9-s3:eq:potencias-Fav-fibonacci} or Cassini’s identity fails;
 \item the rational brackets are not nested or do not contract their width;
 \item the chronological frequency, the uniform measure on the
       catalog or the Parry measure are identified;
 \item a decimal table is involved before the Perron ray is constructed;
 \item a selected cylinder does not truncate to the previous cylinder.
\end{enumerate}
\end{criterion}


% Unidad U015. Acoplamiento de horquillas racionales, refinamiento TPK y
% certificación a profundidad arbitraria.

\section{Diagonal certification, inverse limit and generation to arbitrary precision}
\label{p12-83-c9-s4:chap:precision-arbitraria-constantes}

\subsection{Typed separation of the \(4\) indices of nonadic prolongation}

This unit has the status of subsequent certification. The prefixes and their
compatible extensions come from the internal relation
\(\operatorname{Ext}\), coinductive survival and the APP--TRIT--TPK
transition laws. The rational brackets that follow do not choose
among the \(729\) elements of a fiber: they locate in the Archimedean chart
the already generated extension and certify its precision at arbitrary depth.

Four parameters of different kinds enter each character \(\chi\in\{\pi,e,\varphi\}\):
\begin{enumerate}
 \item \(n\), order of the rational analytic bracket;
 \item \(K\), refinement level of the Topological Phase Kernel;
 \item \(L_K=6K\), ternary length of the published prefix;
 \item \(M\), the requested number of decimal positions.
\end{enumerate}
Neither \(n=K\), \(n=L_K\), \(K=M\) nor a fixed proportionality
between them is postulated. The coupling is constructed through interval inclusions
decided by integer arithmetic.

The integer parts are
\[
 m_\pi=3,
 \qquad
 m_e=2,
 \qquad
 m_\varphi=1,
\]
and the fractional characters are
\[
 \rho_\pi=\pi-3,
 \qquad
 \rho_e=e-2,
 \qquad
 \rho_\varphi=\varphi-1.
\]
The three families of semiopen forks are written
\[
 J_{\chi,n}=[\ell_{\chi,n},h_{\chi,n}).
\]
Their endpoints are rational, contain \(\rho_\chi\) and satisfy
\[
 \ell_{\chi,n}\nearrow\rho_\chi,
 \qquad
 h_{\chi,n}\searrow\rho_\chi,
 \qquad
 h_{\chi,n}-\ell_{\chi,n}\longrightarrow0.
\]
For \(\pi\), the endpoints come from the alternating sums of Machin’s
identity proved in \cref{chap:biografia-pi-autonoma}; for \(e\),
from the factorial sums and their rational remainder in
\cref{p12-83-c9-s2:chap:biografia-completa-e}; for \(\varphi\), from the Fibonacci
and Cassini quotients in \cref{p12-83-c9-s3:chap:biografia-completa-phi}.

\subsection{Ternary cylinders and the immediate ambient fibre}

A ternary prefix of length \(L_K=6K\) is represented by an integer
\(N_K\) with
\[
 0\leq N_K<3^{L_K}.
\]
His semi-open cylinder is
\begin{equation}
 C(N_K,L_K)=
 \left[
 \frac{N_K}{3^{L_K}},
 \frac{N_K+1}{3^{L_K}}
 \right).
 \label{p12-83-c9-s4:eq:cilindro-nivel-K}
\end{equation}
The next emission adds six trits. The immediate ambient fiber contains
\(3^6=729\) elements, indexed by \(u\in\{0,\ldots,728\}\):
\begin{equation}
 C_{K,u}=
 \left[
 \frac{729N_K+u}{3^{L_K+6}},
 \frac{729N_K+u+1}{3^{L_K+6}}
 \right).
 \label{p12-83-c9-s4:eq:729-subcilindros}
\end{equation}
The \(C_{K,u}\) are disjoint and
\[
 C(N_K,L_K)=\bigsqcup_{u=0}^{728}C_{K,u}.
\]
The number \(729\) counts coordinates of this ambient fiber; it does not count
seeds, distinct emissions, orbits, or surviving states.

\subsubsection{Diophantine comparison of the rational endpoints of the gap}

Let \(J=[\ell,h)\subset C(N_K,L_K)\), with
\(\ell,h\in\mathbb Q\). Escribimos \(S_K=3^{L_K+6}\) and define
\begin{align}
 j_{\min}(J,K)
 &=\max\!\left(729N_K,\left\lfloor S_K\ell\right\rfloor\right),
 \label{p12-83-c9-s4:eq:jmin-diagonal}\\
 j_{\max}(J,K)
 &=\min\!\left(729N_K+728,\left\lceil S_Kh\right\rceil-1\right).
 \label{p12-83-c9-s4:eq:jmax-diagonal}
\end{align}

\begin{lemma}[Unitary localization criterion]
\label{p12-83-c9-s4:lem:criterio-localizacion-unitaria}
The bracket \(J\) is contained in a unique element of the fiber
\eqref{p12-83-c9-s4:eq:729-subcilindros} if and only if
\begin{equation}
 \boxed{j_{\min}(J,K)=j_{\max}(J,K).}
 \label{p12-83-c9-s4:eq:jmin-igual-jmax}
\end{equation}
In that case, if \(j\) is the common value,
\[
 u=j-729N_K,
 \qquad
 N_{K+1}=j.
\]
\end{lemma}

\begin{proof}
The integers \(j\) whose elementary intervals of length \(S_K^{-1}\)
intersect \(J\) form the integer interval
\[
 \left[
 \lfloor S_K\ell\rfloor,
 \lceil S_Kh\rceil-1
 \right]\cap
 [729N_K,729N_K+728].
\]
This contains exactly one element if and only if its endpoints coincide.
The equality identifies the numerator of the subcylinder and, by truncating its last
six trits, \(N_K\) is recovered.
\end{proof}

\subsection{Minimal diagonal certification}

Generation does not fix the same analytic order for all
levels in advance. Once the compatible extension has been produced, the order of the
bracket is increased only until it is certified uniquely.

\begin{definition}[Diagonal certification recurrence]
\label{p12-83-c9-s4:def:recurrencia-diagonal-constantes}
Given a family of compatible cylinders fixed by \(\operatorname{Ext}\),
\(C(N_K,L_K)\), whose publication contains \(\rho_\chi\), an initial order
\(n_\chi(K_0-1)\geq1\) is chosen and we recursively define
\begin{equation}
 n_\chi(K)=
 \min\left\{
 n\geq n_\chi(K-1):
 j_{\min}(J_{\chi,n},K)=j_{\max}(J_{\chi,n},K)
 \right\},
 \label{p12-83-c9-s4:eq:nchi-K-minimo}
\end{equation}
and we verify
\begin{equation}
 N_{\chi,K+1}
 =j_{\min}(J_{\chi,n_\chi(K)},K).
 \label{p12-83-c9-s4:eq:Nchi-K-mas-uno}
\end{equation}
The set
\[
 \Delta_\chi=
 \{(n_\chi(K),K):K\geq K_0\}
\]
is the analytic--nonadic diagonal of the character \(\chi\).
\end{definition}

\begin{theorem}[Existence and uniqueness of the certification diagonal]
\label{p12-83-c9-s4:thm:existencia-unicidad-diagonal}
For \(\chi\in\{\pi,e,\varphi\}\), the recurrence
\eqref{p12-83-c9-s4:eq:nchi-K-minimo} is defined at every finite level. The index
\(n_\chi(K)\) and the coordinate \(N_{\chi,K+1}\) are unique.
\end{theorem}

\begin{proof}
Suppose the cylindrical extension of level \(K\) containing
\(\rho_\chi\) has been generated internally. Since \(\rho_\chi\) is irrational, it belongs to no
ternary boundary \(3^{-L}\mathbb Z\). It therefore lies in the interior of a
unique subcylinder \(C_{K,u_*}\), at a positive distance from both endpoints.
The brackets \(J_{\chi,n}\) contain \(\rho_\chi\) and their diameter
tends to zero; hence there exists an \(n\) for which the bracket is
contained in \(C_{K,u_*}\). The set over which the minimum is taken in
\eqref{p12-83-c9-s4:eq:nchi-K-minimo} is nonempty and, by the well-ordering of
\(\mathbb N\), has a least element.
\cref{p12-83-c9-s4:lem:criterio-localizacion-unitaria} proves that the bracket
locates it uniquely and that the located numerator agrees with the one already
generated. Induction on \(K\) completes the argument.
\end{proof}

\begin{corollary}[Absence of decimal parameters]
The diagonal is computed using sums, products, powers, Euclidean divisions,
floors and ceilings of integers. It receives no decimal digit of \(\pi\), \(e\) or
\(\varphi\) as certification data.
\end{corollary}

\subsection{Compatibility with the joint nonadic connection}

Let \(\widehat x_{\chi,K}\) be the enriched state of channel \(\chi\) at
the level \(K\). Its coordinates include the prefix, the cylinder, the phase,
the quotient, the carry, the unbounded vertical memory, the terminal term,
and the provenance register. The residual phase is
\[
 g_0(K)=[K]_9\in\mathbb Z/9\mathbb Z,
\]
while the public label is
\[
 g(K)=1+((K-1)\bmod9)\in\{1,\ldots,9\}.
\]
The vertical memory \(q_{\rm mem}(K)\in\mathbb Z_{\geq0}\) is not reduced
modulo \(12\); only its dodecaphasic projection is
\[
 \beta(K)=q_{\rm mem}(K)\bmod12
\]
does so.

For each \(K\), let
\[
 \Gamma_{9,K}:\widehat X_K\longrightarrow\widehat X_{K+9}
\]
the holonomy of a phase loop. Then
\[
 g_0(K+9)=g_0(K),
 \qquad
 q_{\rm mem}(K+9)=q_{\rm mem}(K)+1,
 \qquad
 \beta(K+9)=\beta(K)+1\pmod{12}.
\]

\begin{theorem}[Phase return with genealogical conservation]
\label{p12-83-c9-s4:thm:retorno-fase-conservacion-genealogica}
Phase equality does not reset the state:
\[
 g_0(K+9)=g_0(K),
 \qquad
 \widehat x_{\chi,K+9}\ne\widehat x_{\chi,K}.
\]
The depths \(27\), \(54\), and \(108\) are the typed compositions
\begin{align*}
 \Gamma_{27,K}
 &=\Gamma_{9,K+18}\circ\Gamma_{9,K+9}\circ\Gamma_{9,K},\\
 \Gamma_{54,K}
 &=\Gamma_{9,K+45}\circ\cdots\circ\Gamma_{9,K},\\
 \Gamma_{108,K}
 &=\Gamma_{9,K+99}\circ\cdots\circ\Gamma_{9,K}.
\end{align*}
In no case are they restarts of the prefix, the carry, the terminal, or the memory.
\end{theorem}

\begin{proof}
The first coordinate is computed in \(\mathbb Z/9\mathbb Z\), whereas
\(q_{\rm mem}\) belongs to an unbounded monoid. Each application of
\(\Gamma_{9,K}\) adds one unit of memory, six blocks per intermediate
level, and the corresponding cocycle composition. Hence the phase
may coincide while the enriched state does not. The depth formulas
display the indices of each map and prevent composition as though all
were endomorphisms of the same space.
\end{proof}

\subsubsection{Compression, expression, and terminal term}

The operational realization of a loop satisfies
\begin{equation}
 U_\Gamma J=JT+\eta,
 \qquad
 J^*\eta=0,
 \qquad
 I-T^*T=\eta^*\eta.
 \label{p12-83-c9-s4:eq:compresion-expresion-precision}
\end{equation}
The visible strict contraction is reserved for
\[
 M_9=\frac59V_9;
\]
is not identified with the compression operator \(T\). Telescoping at
depth \(L\) preserves the terminal term:
\begin{equation}
 S_0=
 \sum_{r=0}^{L-1}M_9^{*r}D_9^*D_9M_9^r
 +M_9^{*L}S_0M_9^L.
 \label{p12-83-c9-s4:eq:telescopia-terminal-precision}
\end{equation}
Omitting the last summand before proving its disappearance would erase
boundary information and would turn phase return into a false
restart.

\subsection{Inverse system of compatible histories}

Let \(\mathcal H_K^\chi\) be the set of admissible enriched histories
of channel \(\chi\) at depth \(K\), and let
\[
 p_{N,K}:\mathcal H_N^\chi\longrightarrow\mathcal H_K^\chi
 \qquad(N\geq K)
\]
the truncation. We define
\[
 \operatorname{Surv}_K^{(N)}(\chi)
 =p_{N,K}(\mathcal H_N^\chi)
\]
and
\begin{equation}
 \operatorname{Surv}_K(\chi)
 =\bigcap_{N\geq K}\operatorname{Surv}_K^{(N)}(\chi).
 \label{p12-83-c9-s4:eq:supervivencia-todas-profundidades}
\end{equation}
The coinductive object is
\begin{equation}
 X_\chi=
 \varprojlim_K
 \bigl(\operatorname{Surv}_K(\chi),p_{K+1,K}\bigr).
 \label{p12-83-c9-s4:eq:limite-inverso-caracteres}
\end{equation}

\begin{theorem}[Section compatible under enhanced extendibility]
\label{p12-83-c9-s4:thm:seccion-compatible-caracteres}
Suppose further that, over each generated numerical subcylinder, the
relation \(\operatorname{Ext}\) is nonempty, single-valued and compatible with the
truncations. Then \(\operatorname{Ext}\) and survival produce a family
\[
 x_\chi=(\widehat x_{\chi,K})_{K\geq K_0}
\]
which satisfies
\[
 p_{K+1,K}(\widehat x_{\chi,K+1})
 =\widehat x_{\chi,K}.
\]
Therefore, \(x_\chi\in X_\chi\).
\end{theorem}

\begin{proof}
The subcylinder of level \(K+1\) has numerator
\(N_{\chi,K+1}=729N_{\chi,K}+u\). Truncation removes its last six
trits and recovers \(N_{\chi,K}\). The other coordinates are obtained through the
functorial laws of carry, memory, terminal and phase. The single-valuedness
hypothesis for \(\operatorname{Ext}\) fixes the subcylinder and the enriched
extension; the local uniqueness of
\cref{p12-83-c9-s4:lem:criterio-localizacion-unitaria} subsequently certifies its
Archimedean localization, and the compatibility of \(\operatorname{Ext}\)
provides truncation equality.
\end{proof}

\subsection{Decimal cells and prescribed precision}

For \(M\geq1\) and \(q\in\{0,\ldots,10^M-1\}\), let
\[
 D_{M,q}=
 \left[
 \frac{q}{10^M},
 \frac{q+1}{10^M}
 \right).
\]
A bracket \(J=[\ell,h)\) is contained in a unique decimal cell if
and only if
\begin{equation}
 \boxed{
 \left\lfloor10^M\ell\right\rfloor
 =\left\lceil10^Mh\right\rceil-1=q.}
 \label{p12-83-c9-s4:eq:criterio-celda-decimal}
\end{equation}

\begin{definition}[Minimum level of decimal certification]
For \(\chi\in\{\pi,e,\varphi\}\), we defines \(K_\chi(M)\) as the least
level for which there exist an analytic order \(n\) and an integer \(q\) such
that
\begin{enumerate}
 \item \(J_{\chi,n}\) satisfies
       \eqref{p12-83-c9-s4:eq:criterio-celda-decimal};
 \item \(J_{\chi,n}\) is contained in the cylinder generated by
       state \(\widehat x_{\chi,K}\);
 \item the transition to the next level satisfies
       \eqref{p12-83-c9-s4:eq:jmin-igual-jmax}.
\end{enumerate}
The pair \((K,n)\) is minimized lexicographically.
\end{definition}

\begin{theorem}[Certification of coinductive outputs to arbitrary decimal precision]
\label{p12-83-c9-s4:thm:precision-decimal-arbitraria}
For each \(\chi\in\{\pi,e,\varphi\}\) and each \(M\geq1\), there exists
\(K_\chi(M)<\infty\). The generated section determines a unique decimal prefix
of \(M\) positions, and the bracket recognizes and certifies that prefix. If
\(M'<M\), the prefix of order \(M\) truncates to the one of
order \(M'\), and the corresponding cells are nested.
\end{theorem}

\begin{proof}
The three characters are irrational: \(\varphi\) is irrational by its quadratic
polynomial with nonsquare discriminant; the irrationality of \(e\) and
\(\pi\) is taken from their classical characterization theorems, subsequent to
the finite construction. Thus none belongs to a rational decimal or
ternary boundary. For a fixed \(M\), the distance from the character to the finite
set of relevant decimal boundaries is positive. The rational
brackets contract their diameter until they lie in a unique decimal cell.
\cref{p12-83-c9-s4:thm:existencia-unicidad-diagonal} couples that bracket to a
TPK cylinder and a unique extension. The lexicographic minimum exists by
well-ordering. Compatibility of the section proves nesting as
\(M\) increases.
\end{proof}

\begin{corollary}[Real uniqueness]
The nested cells have diameters \(10^{-M}\to0\). Their intersection is
a singleton, identified by the preceding characterizations with
\(\pi\), \(e\) or \(\varphi\), respectively.
\end{corollary}

\subsection{Equivalence of \(2\) finite realizations of the uniform law}

The witnesses of \(1,000\) and \(10,000\) positions do not define the
characters or limit the prolongation. They are executions of the same uniform
recurrence at two different depths.

\begin{table}[htbp]
\centering
\caption{Exact dimensions of certified runs, by channel.}
\label{p12-83-c9-s4:tab:dimensiones-ejecuciones-1000-10000}
\begin{tabular}{lrr}
\toprule
Counted object & \(1,000\) positions & \(10,000\) positions\\
\midrule
Partitions exhaustive of the fiber ambient & 396 & 3501\\
Trits transportados & 2376 & 21006\\
Pairs of equal phase \((K,K+9)\) & 387 & 3492\\
Complete nonadic turns & 43 & 388\\
Elements examined by partition & 729 & 729\\
Compatible extensions preserved by partition & 1 & 1\\
\bottomrule
\end{tabular}
\end{table}

Para \(10,000\) posiciones,
\[
 B=3501=389\cdot9,
 \qquad
 6B=21006,
 \qquad
 3^{6B}>10^{10000}.
\]
The first \(396\) blocks of the long run coincide with the
run of \(1,000\) positions. In each of the \(3492\) pairs of
equal phase, the phase coordinate is preserved and the prefix, depth,
cylinder, and memory are extended. The certificate also records any
repetitions of local projections without conflating them with the return of
the complete state.

\subsection{Separation between generation and verification}

\begin{definition}[Structural generator]
The generator receives the finite catalog, \(L_0,L_1\), the initial state, the
connection laws and the internal relation \(\operatorname{Ext}\). At each level, it:
\begin{enumerate}
 \item enumerates the \(729\) elements of the ambient fiber;
 \item applies \(\operatorname{Ext}\) and retains the unique compatible extension;
 \item updates phase, quotient, carry, memory, and terminal;
 \item writes the ternary block, the selection receipt, and the genealogical record.
\end{enumerate}
\end{definition}

\begin{definition}[Independent post hoc verification]
The independent subsequent check receives the closed outputs of the
generator and the three families of rational brackets; it computes
\(j_{\min}\) and \(j_{\max}\), certifies their equality with the extension already
produced, reconstructs the
partitions, checks the truncation squares, the pairs
\((K,K+9)\), the partition sums
\[
 \operatorname{killed}_{\mathrm{izq}}+1+\operatorname{killed}_{\mathrm{der}}=729,
\]
the integrity of the outputs and, finally, the agreement with reference expansions.
\end{definition}

\begin{theorem}[Isolation of the target data]
\label{p12-83-c9-s4:thm:aislamiento-datos-objetivo-precision}
Reference decimal expansions and metrological data are not
involved in selecting a compatible extension. The comparison is
performed after the three prefixes have been written.
\end{theorem}

\begin{proof}
The mathematical input to each generative decision consists of the finite
catalog, \(L_0,L_1\), the inherited state and the internal relation
\(\operatorname{Ext}\). The rational endpoints, powers of \(3\), floors and
ceilings belong to the subsequent check. The generator finishes before
that check opens the brackets or runs the decimal controls. The direction
of the data flow is
\[
 \text{structure}\longrightarrow
 \text{prefixes and ledgers}\longrightarrow
 \text{independent subsequent check};
\]
there is no arrow in the opposite direction.
\end{proof}

\subsection{Differentiated obstructions to generation and its certification at arbitrary precision}

\begin{criterion}[Refutation of generation to arbitrary precision]
Failures of orbit, extension, truncation or return refute the generation;
failures of brackets, order or decimal control refute the corresponding
analytic certificate; and a discrepancy between the two refutes the
joint identification of the published output. With this separation of types,
the complete assertion is refuted if:
\begin{enumerate}
 \item any of the indices \(n,K,L_K,M\) are identified without a
       linking theorem;
 \item a bracket ceases to contain its character or does not contract its diameter;
 \item a stage satisfies \(j_{\min}\ne j_{\max}\) after the order
       declared to be minimal;
 \item an extension does not truncate to the cylinder of the previous level;
 \item the phase return resets the prefix, carry, terminal or memory;
 \item the terminal term in
       \eqref{p12-83-c9-s4:eq:telescopia-terminal-precision} is omitted without proving that it vanishes;
 \item the first \(396\) blocks of the long execution differ from the
       short execution;
 \item a decimal or metrological control is read before the output is closed;
 \item a printed digit differs from the certified mathematical section;
 \item the proof depends on the requested precision being \(1,000\) or
       \(10,000\), rather than an arbitrary \(M\).
\end{enumerate}
\end{criterion}
